% !TEX root = birch2019pressure.tex
% FICTIONAL TEST MATERIAL: this paper does not exist; it was written for the
% citation-audit fixture.
\documentclass{article}
\usepackage{amsmath}

\title{Night-Time Pressure Reduction and Burst Frequency}
\author{C. Fixture \and D. Fixture}
\date{Journal of Fixture Studies, 2019}

\begin{document}
\maketitle

\noindent\textbf{FICTIONAL TEST MATERIAL: this paper does not exist; it was
written for the citation-audit fixture.} No author, utility, venue, or
measurement described here is real.

\begin{abstract}
Water mains burst more often when they run at high pressure, and network
pressure peaks at night, when demand is lowest. We report a two-year trial on
a fictional utility in which pressure-reducing valves lowered night-time
pressure by 20\% in 12 districts. Burst frequency in those districts fell by
28\% relative to the two years before the trial. Because every burst releases
water until it is isolated, the fall in bursts also reduced the water lost to
bursts. We describe the valve settings, the burst records, and the comparison
with untreated districts.
\end{abstract}

\section{Introduction}
\label{sec:introduction}

Each burst on a distribution main costs a repair, an interruption of supply,
and the water that escapes before the main is isolated. Bursts are therefore a
direct source of lost water. Pressure is one cause of bursts that an operator
can control. At night, consumption drops and the pressure in a district rises
towards the static head of its supply. Pressure-reducing valves can cut that
night-time excess while leaving daytime service unchanged.

This paper asks a single question: how much does a sustained cut in night-time
pressure reduce burst frequency? We answer it with a two-year trial on the
fictional Fixture Plains network.

\section{Method}
\label{sec:method}

\subsection{Districts and valve settings}
\label{sec:method-districts}

Twelve hydraulically isolated districts took part. Each is fed through a single
pressure-reducing valve. Before the trial, every valve held a fixed outlet
pressure. During the trial, the valves switched to a time-based profile that
lowered the outlet pressure by 20\% between midnight and five in the morning,
and restored the original setting for the rest of the day. Daytime pressure was
not changed.

\subsection{Burst records}
\label{sec:method-records}

A burst is a main break that required an unplanned repair. We counted bursts
from the utility's repair records for the two years before the trial and for
the two trial years, and normalised the counts by main length. Twelve districts
with similar pipe age and material, whose valves were not changed, served as a
comparison group.

\section{Results}
\label{sec:results}

\subsection{Burst frequency}
\label{sec:results-bursts}

Cutting night-time pressure by 20\% reduced burst frequency by 28\% across the
12 trial districts over the two-year trial, relative to the two years before
it. The comparison districts showed no significant change over the same
period, so the reduction is not explained by a mild winter or by a change in
repair practice. The fall was present in both trial years and in 11 of the 12
districts.

\subsection{Water lost to bursts}
\label{sec:results-losses}

We estimated the water lost to each burst from the repair record and the
recorded time until the main was isolated. Because the trial districts had
fewer bursts, the estimated volume of water lost to bursts fell by 24\% over
the trial. The fall is slightly smaller than the fall in burst frequency,
because the bursts that remained were, on average, on larger mains.

\subsection{Service levels}
\label{sec:results-service}

No customer complaint of low pressure was traced to the night-time profile, and
fire-flow tests at the critical point of each district remained within the
required limits.

\section{Limitations}
\label{sec:limitations}

The trial was designed to measure one effect, and its conclusions are limited
to that effect.

This study does not examine leak localization: we did not attempt to find or
pinpoint leaks, and the trial says nothing about how accurately leaks can be
located. This study does not examine metering: no customer or district meters
were installed or read for the trial, and the loss estimates rest on repair
records rather than on metered flows. The trial ran for two years in twelve
districts; a longer record would be needed to separate the pressure effect from
the slow ageing of the mains.

\end{document}
